\Unit[Hydrostatics (pressure/buoyancy); continuity; Bernoulli; ideal gas; First Law and engines; entropy]{Fluids \& Thermodynamics}
\label{unit:fluids-thermo}

\begin{KnowledgePoint}{Hydrostatics and Buoyancy \Core}
  \KPFormulas
  \begin{formulabox}
  \textbf{Concept explanation:} Static fluids exert pressure that increases with depth; the buoyant force equals the weight of displaced fluid (Archimedes).

  \textbf{Core formulas:}
  \[
  \left\{\begin{aligned}
    &P=P_0+\rho g h\ (\text{hydrostatic pressure}),\\
    &F_b=\rho g V_{\text{disp}}\ (\text{buoyancy}).
  \end{aligned}\right.
  \]

  \textbf{Variable definitions:} $P_0$ reference pressure (often atmospheric at $h=0$); $\rho$ fluid density; $h$ depth; $V_{\text{disp}}$ displaced volume.

  

  \textbf{Prerequisites \& scope:} Fluid at rest (no flow), constant $\rho$ with depth (or integrate if varying); neglect surface tension unless specified.
  \end{formulabox}

  \KPHeuristics
  \begin{heuristicsbox}
  \begin{itemize}[leftmargin=*]
    \item Draw free-body diagrams (FBD) of floating/sinking bodies: set $F_b$ vs weight vs any tension to solve equilibrium.
    \item Choose a definite reference level for $h$ and keep $P_0$ the same when comparing points.
    \item Submerged fraction for floating: $\dfrac{V_{sub}}{V}=\dfrac{\rho_b}{\rho}$ (with $0<\rho_b\le\rho$).
    \item If $\rho_b>\rho$ and the object is released freely, the initial net force is downward: $mg-F_b>0$.
    \item If later supported from below or by upward tension, static equilibrium requires $N+T+F_b=mg$, omitting forces that are absent.
    \item Gauge vs absolute pressure: $\Delta P=\rho g h$ is a \emph{gauge} difference; absolute pressure is $P= P_{\text{atm}}+\rho g h$ when the surface is open to atmosphere.
  \end{itemize}
  \vspace{0.4em}
  \begin{itemize}[leftmargin=*]
    \item Using object's volume instead of displaced volume for $F_b$. Fix: use actual displaced fluid volume (submerged part only).
  \end{itemize}
  \end{heuristicsbox}

  \KPProblems
\begin{cheatproblem}
  A block of volume $V$ and density $\rho_b$ floats in a liquid of density $\rho$. What fraction of its volume is submerged?
  \begin{solutionbox}
  At equilibrium $\rho g V_{sub}=\rho_b g V$, so $V_{sub}/V=\rho_b/\rho$.
  \end{solutionbox}
  \end{cheatproblem}
% (duplicate problem removed)
\end{KnowledgePoint}

\begin{KnowledgePoint}{Continuity and Bernoulli \Core}
  \KPFormulas
  \begin{formulabox}
  \textbf{Concept explanation:} In steady incompressible flow, mass conservation gives $Av=\text{const}$; along a streamline with negligible viscosity, mechanical energy per volume is constant (Bernoulli).

  \textbf{Core formulas:}
  \[
  \left\{\begin{aligned}
    &\text{Continuity: } A_1 v_1=A_2 v_2\ (\rho \ \text{constant}).\\
    &\text{Bernoulli: } P+\tfrac12\rho v^2+\rho g y=\text{const}\ (\text{along a streamline}).
  \end{aligned}\right.
  \]

  \textbf{Variable definitions:} $A$ cross-sectional area; $v$ speed; $P$ pressure; $y$ elevation; $\rho$ density.

  

  \textbf{Prerequisites \& scope:} Steady, incompressible, non-viscous flow; apply Bernoulli along a streamline, not across shocks or with pumps/turbines unaccounted.
  \end{formulabox}

  \KPHeuristics
  \begin{heuristicsbox}
  \begin{itemize}[leftmargin=*]
    \item Along a streamline, check steady, incompressible, inviscid flow before using Bernoulli. An irrotational flow also permits the same constant across streamlines.
    \item Combine continuity with Bernoulli to eliminate speeds or pressures efficiently.
  \item Use stagnation points: where $v=0$, the total (stagnation) pressure is $P_0=P+\tfrac12\rho v^2$ upstream along a streamline.
  \end{itemize}
  \vspace{0.4em}
  \begin{itemize}[leftmargin=*]
    \item Using Bernoulli across different streamlines where viscous losses or pumps exist. Fix: apply along a single streamline and include head gains/losses when needed.
  \end{itemize}
  \end{heuristicsbox}

  \KPProblems
\begin{cheatproblem}
  Water flows through a horizontal pipe from diameter $D_1=0.1\,\text{m}$ to $D_2=0.05\,\text{m}$. If $v_1=2\,\text{m/s}$, find $v_2$ using continuity.
\begin{solutionbox}
  $A_1v_1=A_2v_2\Rightarrow\pi(D_1/2)^2v_1=\pi(D_2/2)^2v_2\Rightarrow v_2=v_1\dfrac{D_1^2}{D_2^2}=2\dfrac{0.01}{0.0025}=8\,\text{m/s}$.
\end{solutionbox}
\end{cheatproblem}
\begin{cheatproblem}
  At point 1 in a pipe, $P_1=1.0\times10^5\,\text{Pa}$, $v_1=2\,\text{m/s}$, $y_1=0$. At point 2, $v_2=5\,\text{m/s}$, $y_2=3\,\text{m}$. Find $P_2$ for water ($\rho=1000\,\text{kg/m}^3$).
\begin{solutionbox}
  Bernoulli: $P_1+\tfrac12\rho v_1^2+\rho gy_1=P_2+\tfrac12\rho v_2^2+\rho gy_2$.
  Compute in SI with scientific notation:
  $\tfrac12\rho(v_1^2-v_2^2)=0.5\times10^3\,(4-25)=-1.05\times10^4\,\text{Pa}$ and
  $\rho g(y_2-y_1)=10^3\times9.8\times3=2.94\times10^4\,\text{Pa}$.
  Hence $P_2=1.00\times10^5-1.05\times10^4-2.94\times10^4=6.01\times10^4\,\text{Pa}$.
\end{solutionbox}
\end{cheatproblem}
\end{KnowledgePoint}

\begin{KnowledgePoint}{Ideal Gas and First Law \Core}
  \KPFormulas
  \begin{formulabox}
  \textbf{Concept explanation:} $PV=nRT$ holds for ideal gases; the First Law links changes in internal energy to heat and work with clear sign conventions.

  \textbf{Core formulas:}
  \[
  \left\{\begin{aligned}
      &PV=nRT,\quad U=\tfrac{f}{2}nRT,\\
      &\text{where } f\ \text{is dof (e.g., }f=3\ \text{monatomic},\ 5\ \text{diatomic at room T)},\\
    &\Delta U=Q+W_{\text{on}}\ \ (W_{\text{on}}=\text{work done on the system}),\\
    &W_{\text{on}}=-\int P_{\rm ext}\,dV,\quad W_{\text{by}}=-W_{\text{on}}.
  \end{aligned}\right.
  \]

  \textbf{Variable definitions:} $P,V,T$ pressure/volume/temperature; $n$ moles; $R$ gas constant; $Q$ heat into system; $W_{\text{on}}$ work on system; $W_{\text{by}}$ work done by the gas; $U$ internal energy.

  

  \textbf{Prerequisites \& scope:} Ideal gas approximation; $U$ depends only on $T$ for ideal gases. The stated $f$ values use classical equipartition. For quasistatic expansion $P_{\rm ext}=P$; in general use the external pressure for boundary work.
  \end{formulabox}

  \KPHeuristics
  \begin{heuristicsbox}
  \begin{itemize}[leftmargin=*]
    \item Identify process (isochoric/isobaric/isothermal/adiabatic) to pick $W,\,Q,\,\Delta U$ quickly.
    \item Draw $P$–$V$ diagrams: areas give work; direction indicates sign.
  \end{itemize}
  \vspace{0.4em}
  \begin{itemize}[leftmargin=*]
    \item Mixing sign conventions for work. Fix: adopt $\Delta U=Q+W_{\text{on}}$ (work on system positive) consistently; then $W_{\text{by}}=-W_{\text{on}}$.
  \end{itemize}
  \end{heuristicsbox}

  \KPProblems
\begin{cheatproblem}
  An ideal gas undergoes an isothermal expansion from volume $V_1$ to $V_2$ at temperature $T$. Find the work done by the gas.
  \begin{solutionbox}
  $W=\int_{V_1}^{V_2} \frac{nRT}{V}\,dV=nRT\ln\!\frac{V_2}{V_1}$.
  \end{solutionbox}
  \end{cheatproblem}
\begin{cheatproblem}
  An ideal gas at $P=2\times10^5\,\text{Pa}$ and $V=0.01\,\text{m}^3$ has $n=1\,\text{mol}$. Find the temperature $T$ using $PV=nRT$ with $R=8.314\,\text{J/(mol·K)}$.
\begin{solutionbox}
  $T=\dfrac{PV}{nR}=\dfrac{2\times10^5\cdot0.01}{1\cdot8.314}\approx240\,\text{K}$.
  \end{solutionbox}
  \end{cheatproblem}
\end{KnowledgePoint}

\begin{KnowledgePoint}{Entropy and Carnot \Core}
  \KPFormulas
  \begin{formulabox}
  \textbf{Concept explanation:} The total entropy of an isolated system increases in an irreversible process; the entropy of a subsystem can decrease. Carnot gives the maximum efficiency of engines operating between two reservoirs.

  \textbf{Core formulas:}
  \[
  \left\{\begin{aligned}
    &\Delta S=\int_{\text{rev}} \frac{\delta Q}{T},\\
    &\text{Carnot bound: } \eta_{\max}=1-\frac{T_c}{T_h}.
  \end{aligned}\right.
  \]

  \textbf{Variable definitions:} $S$ entropy; $T_h,T_c$ hot/cold absolute temperatures; $\delta Q$ infinitesimal heat (reversible path).

  

  \textbf{Prerequisites \& scope:} Absolute temperatures (Kelvin); reversible paths for definition; real engines achieve less than Carnot due to irreversibilities.
  \end{formulabox}

  \KPHeuristics
  \begin{heuristicsbox}
  \begin{itemize}[leftmargin=*]
    \item Compute $\Delta S$ along a convenient reversible path (e.g., isothermal + isochoric steps).
    \item For engine limits, compare cycle temperatures to $T_h,T_c$ to bound $\eta$ quickly.
  \end{itemize}
  \vspace{0.4em}
  \begin{itemize}[leftmargin=*]
    \item Using Celsius in $\eta=1-T_c/T_h$. Fix: convert to Kelvin.
  \end{itemize}
  \end{heuristicsbox}

  \KPProblems
\begin{cheatproblem}
  A Carnot engine operates between $T_h=500\,\text{K}$ and $T_c=300\,\text{K}$. Find its maximum efficiency.
\begin{solutionbox}
  $\eta_{\max}=1-\dfrac{T_c}{T_h}=1-\dfrac{300}{500}=0.4=40\%$.
\end{solutionbox}
\end{cheatproblem}
\begin{cheatproblem}
  An ideal gas expands reversibly at constant $T=400\,\text{K}$ from $V_1=1\,\text{m}^3$ to $V_2=2\,\text{m}^3$. Find the entropy change if $n=1\,\text{mol}$.
\begin{solutionbox}
  $\Delta S=\dfrac{Q_{rev}}{T}=\dfrac{nRT\ln(V_2/V_1)}{T}=nR\ln2=8.314\ln2\approx5.76\,\text{J/K}$.
\end{solutionbox}
\end{cheatproblem}
\end{KnowledgePoint}

\begin{KnowledgePoint}{Heat Engines and Efficiency \Core}
  \KPFormulas
  \begin{formulabox}
  \textbf{Core formulas:}
  \[
  \left\{\begin{aligned}
    &W_{\text{by}}=\int P_{\rm ext}\,dV,\quad W_{\text{on}}=-W_{\text{by}},\quad \eta=\frac{W_{\text{net,by}}}{Q_h}.
  \end{aligned}\right.
  \]
  \textbf{Variable definitions:} $Q_h$ is heat absorbed from the hot reservoir; $W_{\text{net,by}}$ is net work per cycle. The Carnot bound is given in the preceding section. For quasistatic steps, $P_{\rm ext}=P$.
  \end{formulabox}
\end{KnowledgePoint}


\UnitProblemLinks{fluids-thermo}
